\documentclass[12pt,a4paper,finnish]{tutthesis}
 
% LaTeX-tiedosto opinnäytepohjaksi
% Tekijä: Sami Paavilainen
% Muokkaaja: Heikki Huttunen (heikki.huttunen@tut.fi) 31.7.2012.
 
% Vaatii lisäksi luokkatiedoston tutthesis.cls
% sekä tiedoston tty-logo.pdf (pdflatex) tai tty-logo-eps (latex)
 
\usepackage[utf8]{inputenc}








\usepackage[small]{caption}% kuvatekstien kirjasinkoko ja asettelu
%omia käskyjä voi lisätä tähän väliin, esim. \newcommand{\angs}{\textsl{\AA}}
\pagenumbering{Roman}
% alkuvalmistelut loppuu
 
\begin{document}
 
\thispagestyle{empty}
 
\vspace*{-.5cm}\noindent
 
% If thesis is in English, use the file "tut-logo"
% instead of "tty-logo" in the following:
 
\includegraphics[width=8cm]{tut-logo}
 
\vspace{6.8cm}
 
\noindent{\bf\large \textsf{Mikko Pohja}}\\
{\bf\large \textsf{Maximizing quality in a small budget software project}}\\
\textsf{Diplomityö}
 
\vspace{8.7cm} % jos kaksi otsikkoriviä vaihda -> 6.7cm
 
\begin{flushright}
  
\begin{minipage}[c]{6.8cm}
\begin{spacing}{1.0}
\textsf{Tarkastaja: Tarkastaja 1}\\
\textsf{Tarkastaja ja aihe hyväksytty}\\ 
\textsf{xxxxxxx tiedekuntaneuvoston}\\
\textsf{kokouksessa xx.xx.xxxx}\\
\end{spacing}
\end{minipage}
\end{flushright}
 
\newpage
 
\setcounter{page}{1} % tämä tarvitaan, jottei ensimmäinen sivu kansilehden jälkeen olisi numero 2.
 
\chapter*{TIIVISTELMÄ}
\begin{spacing}{1.0}
\textsf{TAMPEREEN TEKNILLINEN YLIOPISTO}\\
\textsf{xxxxxxxxxxxxx koulutusohjelma}\\
{\bf \textsf{Mikko Pohja: Maximizing quality in a small budget software project}}\\
\textsf{Diplomityö, xx sivua, x liitesivua}\\
\textsf{Xxxxxkuu 201x}\\
\textsf{Pääaine: }\\
\textsf{Tarkastajat: }\\
\textsf{Avainsanat: }\\
\end{spacing}
 
\noindent
Ensimmäinen kappale
 
\noindent
Toinen kappale
\newpage
\chapter*{ABSTRACT}
\begin{spacing}{1.0}
\textsf{TAMPERE UNIVERSITY OF TECHNOLOGY}\\
\textsf{Master's Degree Programme in xxxxxxx Technology}\\
{\bf \textsf{AUTHOR : Title}}\\
\textsf{Master of Science Thesis, xx pages, x Appendix pages}\\
\textsf{xxxxxx 201x}\\
\textsf{Major: }\\
\textsf{Examiner: }\\
\textsf{Keywords: }\\
\end{spacing}
 
\noindent 
First paragraph
 
\noindent
Second paragraph
 
\newpage
 
\chapter*{ALKUSANAT}
\noindent Tämä (\textit{d-tyo.tex}) on \LaTeX-pohja Tampereen teknillisen
yliopiston opinnäytetöitä varten. Samaan pakettiin kuuluu myös
tiedosto \mbox{\textit{tutthesis.cls}}, joka sisältää taittoteknisiä
lisäyksiä \LaTeX:n alkuperäiseen \textit{report.cls}-luokkatiedostoon.
 
Lisäksi otsikkosivua varten tarvitaan tiedosto \textit{tut-logo.xxx}, jonka
tulee sisältää TTY:n logo. Tiedoston tulee
olla joko \textit{.eps}- tai \textit{.pdf}-muodossa riippuen \LaTeX-versiosta.
 
\newpage
\tableofcontents
%\newpage
%\listoffigures
%\listoftables
\newpage
\chapter*{TERMIT JA NIIDEN MÄÄRITELMÄT}
 
 
\begin{termiluettelo}
 
\item [$\hbar$] Redusoitu Planckin vakio
\item [SNR] Signaali-kohinasuhde (engl.: Signal to Noise Ratio)
 
\end{termiluettelo} 
 
 
\newpage
\renewcommand{\chaptermark}[1]{\markboth{\thechapter. \ #1}{}}
\renewcommand{\sectionmark}[1]{\markright{}{}}
\lhead{\fancyplain{}{\leftmark}}
 
% Tästä alkaa varsinainen teksti
\pagenumbering{arabic}
 
\chapter{Johdanto}
 
 ECO: Economics of software quality
 ROI: Improving the ROI of software quality--
 
 \section{TODO}




 \begin{itemize}




 \item Technical debt ECO: Chapter 7
 
 \end{itemize} 
 
 \chapter{Quality Assurance in Futurice}
 
TODO

 \chapter{Startup}
 
 \section{Budget}




 \begin{itemize}
  \item Startup financing cycle (Valley of Death)
 \end{itemize}
 
 \section{Life cycle}




 \begin{itemize}
 
 \item Scope
 \item MVP
 \item Oikeat feature
 \item Validoinnit: I Know I When I See It
 \item Lyhyet sprintit
 
 \end{itemize}
 
 \section{Analytics}
 
 \begin{itemize}
 
 \item Ominaisuuksien priorisointi (esim. käytön mukaan)
 
 \end{itemize}
  
 \chapter{Software Quality}
 
 \section{Software quality in General}
 
 \begin{itemize}
 
 \item mittaus
 
 \item tavoitteena luoda arvoa asiakkaalle
 
 \item Motivaatio:
 
 \begin{itemize}
 
 \item ECO: it is gratifying to observe that high quality levels are invariably associated with shorter-than-average development schedules and lower-than-average development costs
 \item ECO: Shortens development schedules, Lowers development costs, Lowers maintenance costs, Reduces warranty costs, Increases customer satisfaction
 \item ECO: Chapter 1, software has become..
 
 \end{itemize}
 
 \item Qualityn määritelmä




 \begin{itemize}
 
 \item ECO: Chapter 1
 
 \end{itemize}
 
 \end{itemize}
 
 \section{Methods for improving quality}
 
 \begin{itemize}
 
 \item ECO: pretest & tests
 
 \item ECO: Chapter 3 preventing defects 
 
 \item ROI: Three main activities: Review, process audit and testing
 
 \end{itemize}
 
 \subsection{Testing}
 
 \begin{itemize}
 
 \item Crash, Smoke and Kattava testaus
 
 \item ECO: Chapter 5
 
 \end{itemize}
 
 \subsection{Reviews, Inspections and other pretest procedures}
 
 \begin{itemize}
  
 \item ECO: chapter 4, pretest defect removal
 
 \end{itemize}
 
 \subsection{Continuous integration}
 
 \subsection{Lean}
 
 \section{Evaluation of QA methods}
 
 \begin{itemize}
  
 \item ECO: We use these quality metrics to compare a number of quality improvement techniques at each stage of the software development life cycle and quantify their efficacy using data from real-world applications.
 
 \end{itemize}
 
 \subsection{Methods in different phases of the life cycle}
 
 \subsection{Point of diminishing returns}
 
 \subsection{The Don't}
 
 \begin{itemize}
 
 \item Liittyy vahvasti Leanin Wasteen
 \item Älä raportoi bugeja, joista tiedät, ettei niitä korjata
 \item Ei turhia raportteja
 \item ECO: Cost per defect => paras tulos bugisimmassa projektissa
 
 \end{itemize]
 
 \chapter{Case: Päikky}
 
 \section{QA methods and processes used}
 
 \begin{itemize}
 
 \item Käytetyt QA-metodit
 \item Kerrotaan prosessista ja tiimistä
 
 \end{itemize}
 
 \section{Achieved quality in the project}




 \begin{itemize}
 
 \item Kartoitetaan bugit (Mailit, Pivotal, Repo)  
 \item Asiakkaan ja loppukäyttäjien tyytyväisyys
 
 \end{itemize}
 
 \chapter{Results}
 
% lähdeluettelo
\newpage
\addcontentsline{toc}{chapter}{Lähteet} % Lisää tämäkin sisällysluetteloon
\renewcommand{\bibname}{Lähteet} 
\begin{thebibliography}{99}
\bibitem{Hirs} Hirsjärvi, S., Remes, P., ja Sajavaara, P. 2005. Tutki ja
  kirjoita, 11. painos. Helsinki, Tammi. 436 s. 
\bibitem{Mittelbach} Mittelbach, F., Goossens, M., Braams, J., 
Carlisle, D., Rowley, C. 2004. The Latex Companion, 2. painos. Boston, Addison-Wesley. 1120 s.
 
\end{thebibliography}
 
% seuraavia voi käyttää bibtex:n kanssa (edellyttää tiedoston kirjat.bib)
%\bibliographystyle{plain}
%\bibliography{kirjat}
 
\appendix
\renewcommand{\chaptername}{Liite}
 
\chapter{Liitteitä}
\end{document}